\textbf{Phần 1. Trắc nghiệm lựa chọn}
\setcounterpageref{ex}{0}

\begin{ex} %[2H14N72]
Hình nào sau đây có 3 cạnh và 3 đỉnh?
\choice
{Hình vuông}
{Hình tứ giác}
{\True Hình tam giác}
{Hình chữ nhật}
\loigiai{
Hình tam giác có 3 cạnh và 3 đỉnh.\\
\textbf{Đáp án của câu hỏi là C}
}
\end{ex}

\begin{ex} %[2H14N72]
Hình tứ giác có bao nhiêu cạnh?
\choice
{$3$ cạnh}
{$5$ cạnh}
{$6$ cạnh}
{\True $4$ cạnh}
\loigiai{
Hình tứ giác có 4 cạnh.\\
\textbf{Đáp án của câu hỏi là D}
}
\end{ex}

\begin{ex} %[2H14H72]
Đường gấp khúc $ABCD$ gồm mấy đoạn thẳng?
\choice
{$2$ đoạn thẳng}
{$4$ đoạn thẳng}
{\True $3$ đoạn thẳng ($AB, BC, CD$)}
{$5$ đoạn thẳng}
\loigiai{
Gồm 3 đoạn thẳng: $AB, BC, CD$.\\
\textbf{Đáp án của câu hỏi là C}
}
\end{ex}

\begin{ex} %[2H14H72]
Độ dài đường gấp khúc gồm hai đoạn thẳng dài $15\text{ cm}$ và $25\text{ cm}$ là
\choice
{\True $40\text{ cm}$}
{$35\text{ cm}$}
{$45\text{ cm}$}
{$50\text{ cm}$}
\loigiai{
$15 + 25 = 40\text{ cm}$.\\
\textbf{Đáp án của câu hỏi là A}
}
\end{ex}

\begin{ex} %[2H14H72]
Khối nào sau đây lăn được về mọi phía?
\choice
{Khối trụ}
{Khối lập phương}
{Khối hộp chữ nhật}
{\True Khối cầu}
\loigiai{
Khối cầu lăn được về mọi phía.\\
\textbf{Đáp án của câu hỏi là D}
}
\end{ex}

\begin{ex} %[2H14H72]
Khối lập phương có bao nhiêu mặt?
\choice
{$4$ mặt}
{$8$ mặt}
{\True $6$ mặt vuông bằng nhau}
{$12$ mặt}
\loigiai{
Khối lập phương có 6 mặt vuông bằng nhau.\\
\textbf{Đáp án của câu hỏi là C}
}
\end{ex}

\begin{ex} %[2H14V72]
Đường gấp khúc gồm 3 đoạn thẳng có độ dài lần lượt là $12\text{ cm}, 18\text{ cm}, 20\text{ cm}$. Độ dài đường gấp khúc đó là
\choice
{$40\text{ cm}$}
{\True $50\text{ cm}$}
{$60\text{ cm}$}
{$45\text{ cm}$}
\loigiai{
$12 + 18 + 20 = 50\text{ cm}$.\\
\textbf{Đáp án của câu hỏi là B}
}
\end{ex}

\begin{ex} %[2H14V72]
Một hình tam giác có độ dài các cạnh là $15\text{ cm}, 20\text{ cm}, 25\text{ cm}$. Chu vi hình tam giác đó là
\choice
{$50\text{ cm}$}
{$70\text{ cm}$}
{$55\text{ cm}$}
{\True $60\text{ cm}$}
\loigiai{
$15 + 20 + 25 = 60\text{ cm}$.\\
\textbf{Đáp án của câu hỏi là D}
}
\end{ex}

\begin{ex} %[2H14V72]
Một hình tứ giác có độ dài 4 cạnh đều bằng $5\text{ cm}$. Chu vi của hình tứ giác đó là
\choice
{$15\text{ cm}$}
{\True $20\text{ cm}$}
{$25\text{ cm}$}
{$10\text{ cm}$}
\loigiai{
$5 \times 4 = 20\text{ cm}$.\\
\textbf{Đáp án của câu hỏi là B}
}
\end{ex}

\begin{ex} %[2H14V72]
Trong các vật sau: hộp sữa chua hình trụ, quả bóng bàn, viên gạch, quả cam. Có bao nhiêu vật có dạng khối cầu?
\choice
{\True $2$ vật (quả bóng bàn, quả cam)}
{$1$ vật}
{$3$ vật}
{$4$ vật}
\loigiai{
Có 2 vật: quả bóng bàn và quả cam.\\
\textbf{Đáp án của câu hỏi là A}
}
\end{ex}

\begin{ex} %[2H14V72]
Một đường gấp khúc gồm 4 đoạn thẳng bằng nhau, mỗi đoạn dài $5\text{ dm}$. Độ dài đường gấp khúc theo mét là
\choice
{$1\text{ m}$}
{\True $2\text{ m}$ ($20\text{ dm}$)}
{$3\text{ m}$}
{$4\text{ m}$}
\loigiai{
$5 \times 4 = 20\text{ dm} = 2\text{ m}$.\\
\textbf{Đáp án của câu hỏi là B}
}
\end{ex}

\begin{ex} %[2H14C72]
Một sợi dây đồng uốn thành một hình tam giác có 3 cạnh bằng nhau, mỗi cạnh dài $8\text{ cm}$. Nếu duỗi thẳng sợi dây rồi uốn thành hình tứ giác có 4 cạnh bằng nhau thì mỗi cạnh của hình tứ giác dài là
\choice
{\True $6\text{ cm}$}
{$5\text{ cm}$}
{$7\text{ cm}$}
{$8\text{ cm}$}
\loigiai{
Độ dài dây: $8 \times 3 = 24\text{ cm}$.\\
Cạnh hình tứ giác: $24 : 4 = 6\text{ cm}$.\\
\textbf{Đáp án của câu hỏi là A}
}
\end{ex}

\textbf{Phần 2. Trắc nghiệm đúng sai}
\setcounterpageref{ex}{0}

\begin{ex} %[2H14H72]
Xét tính đúng sai về các hình hình học:
\choiceTF
{\True Ba điểm thẳng hàng là ba điểm cùng nằm trên một đường thẳng.}
{\True Đường gấp khúc $MNPQ$ gồm 3 đoạn thẳng $MN, NP, PQ$.}
{Khối trụ có hai mặt đáy là hai hình vuông.}
{\True Khối trụ có hai mặt đáy là hai hình tròn bằng nhau.}
\loigiai{
\begin{itemize}
    \item 3 điểm thẳng hàng: Đúng.
    \item Đường gấp khúc MNPQ có 3 đoạn: Đúng.
    \item Đáy khối trụ là hình tròn (không phải hình vuông). Ý 3 Sai, Ý 4 Đúng.
\end{itemize}
\begin{itemize}
    \item a) \textbf{Đúng}.
    \item b) \textbf{Đúng}.
    \item c) \textbf{Sai}.
    \item d) \textbf{Đúng}.
\end{itemize}
}
\end{ex}

\begin{ex} %[2H14V72]
Một hình tứ giác có độ dài các cạnh lần lượt là $10\text{ cm}, 12\text{ cm}, 15\text{ cm}, 13\text{ cm}$. Xét tính đúng sai:
\choiceTF
{\True Chu vi hình tứ giác là $50\text{ cm}$ ($10 + 12 + 15 + 13 = 50$).}
{\True Chu vi hình tứ giác đổi ra đề-xi-mét là $5\text{ dm}$.}
{Chu vi hình tứ giác là $60\text{ cm}$.}
{\True $50\text{ cm} = 5\text{ dm} \neq 60\text{ cm}$.}
\loigiai{
\begin{itemize}
    \item Chu vi = 50 cm: Đúng.
    \item $50\text{ cm} = 5\text{ dm}$: Đúng.
    \item $50 \neq 60$. Ý 3 Sai, Ý 4 Đúng.
\end{itemize}
\begin{itemize}
    \item a) \textbf{Đúng}.
    \item b) \textbf{Đúng}.
    \item c) \textbf{Sai}.
    \item d) \textbf{Đúng}.
\end{itemize}
}
\end{ex}

\textbf{Phần 3. Trả lời ngắn}
\setcounterpageref{ex}{0}

\begin{ex} %[2H14H72]
Hình tứ giác có mấy cạnh?
\shortans{4}
\loigiai{
\textbf{Đáp số:} $4$\\
4 cạnh.
}
\end{ex}

\begin{ex} %[2H14H72]
Khối lập phương có mấy mặt?
\shortans{6}
\loigiai{
\textbf{Đáp số:} $6$\\
6 mặt.
}
\end{ex}

\begin{ex} %[2H14V72]
Đường gấp khúc gồm 2 đoạn thẳng dài $18\text{ cm}$ và $22\text{ cm}$. Tính độ dài đường gấp khúc theo xăng-ti-mét.
\shortans{40}
\loigiai{
\textbf{Đáp số:} $40$\\
$18 + 22 = 40\text{ cm}$.
}
\end{ex}

\begin{ex} %[2H14V72]
Một hình tam giác có 3 cạnh đều dài $7\text{ cm}$. Tính chu vi hình tam giác đó theo xăng-ti-mét.
\shortans{21}
\loigiai{
\textbf{Đáp số:} $21$\\
$7 \times 3 = 21\text{ cm}$.
}
\end{ex}

\begin{ex} %[2H14V72]
Có 8 khối cầu và 7 khối trụ. Hỏi có tất cả bao nhiêu khối hình?
\shortans{15}
\loigiai{
\textbf{Đáp số:} $15$\\
$8 + 7 = 15$ khối.
}
\end{ex}

\begin{ex} %[2H14C72]
Đường gấp khúc gồm 5 đoạn thẳng bằng nhau có tổng độ dài là $35\text{ cm}$. Hỏi mỗi đoạn thẳng dài bao nhiêu xăng-ti-mét?
\shortans{7}
\loigiai{
\textbf{Đáp số:} $7$\\
$35 : 5 = 7\text{ cm}$.
}
\end{ex}
